\documentclass[11pt]{article}
\usepackage{graphicx}
\usepackage{amssymb}
\usepackage{bm}
\usepackage{mathtools}
\usepackage{amsmath}
\usepackage{commath}
\usepackage{amsthm}
\usepackage{enumitem}
\usepackage{float} 
\usepackage{array}
\usepackage{outlines}
\usepackage{setspace}
\usepackage[colorinlistoftodos]{todonotes}
\definecolor{green}{rgb}{0.4, 0.69, 0.2}
\definecolor{lavender(floral)}{rgb}{0.71, 0.49, 0.86}
\definecolor{asparagus}{rgb}{0.53, 0.66, 0.42}
\definecolor{blue}{rgb}{0.0, 0.53, 0.74}
\setlist[enumerate,1]{label=\roman*)}
\setlist[enumerate,2]{label=\alph*)}
\newtheorem{theorem}{Theorem}
\newtheorem{corollary}{Corollary}
\newtheorem{lemma}{Lemma}
\newtheorem{remark}{Remark}
\newtheorem{definition}{Definition}
\newtheorem{prop}{Proposition}
\newcommand*\diff{\mathop{}\!\mathrm{d}}
\usepackage{tikz}
\usetikzlibrary{arrows,shapes,trees, calc}


%%% *** PREAMBLE *** %%%
\title{Strong convergence of strong solutions for the case of $H^3$ initial data ${\mathbf{u}_h}_0$ and downwind-matching coordinate system}

\date{}

%%%***  CONTENT *** %%%
\begin{document}

\maketitle

\begin{abstract}

The strong convergence of $(\mathbf{u}_h^\epsilon, u_3^\epsilon, c^\epsilon)$ is proposed over the unbounded $[0, \infty)$ domain, with $\epsilon$ convergence rate and a non-uniform bound in time.

\end{abstract}

\newpage
\tableofcontents
\newpage

\section{Downwind-matching coordinates}

Boundary conditions: we use the minimal additional set of restrictions that allow us to benefit from a simpler structure that resembles periodic domains, keeping a classical non-periodic domain however. (details in the next section)

\textbf{Change of notations:} Since $u^\epsilon$ and $u$ denotes the solution functions and the capital $U^\epsilon$ denotes their difference, we switch to the analogous notation with the concentration function, i.e. $c^\epsilon$ and $c$ here are the solutions in the anisotropic and hydrostatic equations, respectively, while the capital $C^\epsilon$ denotes the $c^\epsilon - c$ difference function. $C$ stands for a generic constant.

  \subsection{The set of equations}

(E-SNS): extended scaled Navier-Stokes.

    \begin{equation} \label{ANSC1-dw}
    \partial_t u^\epsilon _1 + \mathbf{u}^\epsilon \cdot \nabla u^\epsilon _1 -\Delta_\nu u^\epsilon _1 -\gamma u_2^\epsilon + \epsilon \beta u_3^\epsilon + \partial_1 p^\epsilon = 0 \text{ in } \Omega \times (0,T)
    \end{equation}
    
    \begin{equation} \label{ANSC2-dw}
    \partial_t u_2^\epsilon + \mathbf{u}^\epsilon \cdot \nabla u_2^\epsilon -\Delta_\nu u_2^\epsilon + \gamma u_1^\epsilon  - \alpha \epsilon u^\epsilon_3 + \partial_2 p^\epsilon = 0 \text{ in } \Omega \times (0,T)
    \end{equation}
    
    \begin{equation} \label{ANSC3-dw}
    \epsilon^2 \{  \partial_t u_3^\epsilon + \mathbf{u}^\epsilon \cdot \nabla u_3^\epsilon -\Delta_\nu u_3^{\epsilon} \} + \epsilon \alpha u^\epsilon_2 - \epsilon \beta u_1^\epsilon + \partial_3 p^\epsilon = 0 \text{ in } \Omega \times (0,T)
    \end{equation}
    
    \begin{equation} \label{ANSC4-dw}
    \nabla \cdot \mathbf{u}^\epsilon= 0 \text{ in } \Omega \times (0,T)
    \end{equation}
    
    \begin{equation} \label{ANSC5-dw}
    \partial_t c^\epsilon + \mathbf{u^\epsilon} \cdot \nabla c^\epsilon = \epsilon K_1 \partial_{11} ^ 2 c^\epsilon + K_2 \partial_{22} ^ 2 c^\epsilon + K_3 \partial_{33} ^ 2 c^\epsilon + s^\epsilon \text{ in } \Omega \times (0,T).
    \end{equation}
    

(E-PEs): extended Primitive Equations.

    \begin{equation} \label{HNSC1-dw}    
    \partial_t u_1 + \mathbf{u} \cdot \nabla u_1 - \Delta_\nu u_1 - \gamma u_2 + \partial_1 p = 0 \text{ in } \Omega \times (0,T)
    \end{equation}
    
    \begin{equation} \label{HNSC2-dw}   
    \partial_t u_2 + \mathbf{u} \cdot \nabla u_2 - \Delta_\nu u_2 + \gamma u_1 + \partial_2 p = 0 \text{ in } \Omega \times (0,T)
    \end{equation}
    
    \begin{equation} \label{HNSC3-dw}   
    \partial_3 p = 0 \text{ in } \Omega \times (0,T)
    \end{equation}
    
    \begin{equation} \label{HNSC4-dw}   
    \nabla \cdot \mathbf{u} = 0 \text{ in } \Omega \times (0,T)
    \end{equation}
    
    \begin{equation} \label{HNSC5-dw}   
    \partial_t c + \mathbf{u} \cdot \nabla c = K_2 \partial_{22} ^ 2 c + K_3 \partial_{33} ^2 c + s \text{ in } \Omega \times (0,T)
    \end{equation}
    
\subsection{Governing equation for $C^\epsilon$}

Denote $(U^\epsilon, C^\epsilon) = (V^\epsilon, W^\epsilon, C^\epsilon) = (u^\epsilon - u, c^\epsilon - c) = (v^\epsilon - v, w^\epsilon - w, c^\epsilon - c).$

The equation describing $C^\epsilon:$ calculating the difference between (\ref{ANSC5-dw}) and (\ref{HNSC5-dw}) we have

$$ \partial_t C^\epsilon + u^\epsilon \cdot \nabla c^\epsilon - u \cdot \nabla c = \epsilon K_1 \partial_{11} c^\epsilon + K_2 \partial_{22} C^\epsilon + K_3 \partial_{33} C^\epsilon + S^\epsilon.$$

Rewriting the term $u^\epsilon \cdot \nabla c^\epsilon - u \cdot \nabla c$ we use

\begin{equation}
  \begin{split}
    & u^\epsilon \cdot \nabla c^\epsilon - u \cdot \nabla c + 0 = u^\epsilon \cdot \nabla c^\epsilon - u \cdot \nabla c - u \nabla c^\epsilon + u \nabla c^\epsilon \\
    & = u^\epsilon \cdot \nabla c^\epsilon - u \nabla c^\epsilon - u \cdot \nabla c + u \nabla c^\epsilon = U^\epsilon \nabla c^\epsilon + u \nabla C^\epsilon \\
    & = U^\epsilon \nabla (c^\epsilon - c + c) + u \nabla C^\epsilon = U^\epsilon \nabla C^\epsilon  + U^\epsilon \nabla c+ u \nabla C^\epsilon.
  \end{split}
\end{equation}

Eventually we have

$$ \partial_t C^\epsilon + U^\epsilon \nabla C^\epsilon  + U^\epsilon \nabla c+ u \nabla C^\epsilon = \textcolor{orange}{\epsilon K_1 \partial_{11} C^\epsilon + \epsilon K_1 \partial_{11} c} + K_2 \partial_{22} C^\epsilon + K_3 \partial_{33} C^\epsilon + S^\epsilon. $$

Multiplying the equation by $- \Delta C^\epsilon$ results only $\epsilon \norm{\partial_{11} C^\epsilon}$ on the left-hand side, while from the estimates of the convective terms, we still have an $\epsilon$-independent $\zeta \norm{\Delta V^\epsilon}^2_2$ term on the right-hand side of the estimate, and these can not anymore be "melted" into the left-hand side, as the full, $\epsilon$-free Laplacian is missing. We will take the following different strategy where the Laplacians are not all brought to the left-hand side.

\subsection{Description}

Let's have ${\mathbf{u}_h}_0 \in H^3(\Omega)$ (and as a consequence $\mathbf{u}_0 \in H^2(\Omega)$) and $\epsilon < \epsilon_0,$ where $\epsilon_0$ is the limit specified by Proposition 5.3 in the Li-Titi article. Then we have the following cornerstones needed for estimating the difference function $C^\epsilon.$

\textbf{Cornerstone existence results and regularities:}

\begin{enumerate}

  \item \textbf{Existence for (E-PEs).} Global unique existence of the $(\mathbf{u}_h, u_3, c)$ strong solution of the (E-PEs) expressed in downwind-matching coordinates: in Cao-Li-Titi's paper \todo{remains to be checked} it has been proved, a difference is that we do not have the $c$-analogous version of the $\partial_z p$---$T$ relationship.

  \item \textbf{Existence for (E-SNS).} From the Li-Titi article we have the global (considering that we are at the case of $\epsilon < \epsilon_0$ as previously specified) unique existence of the $(\mathbf{u}_h^\epsilon, u_3^\epsilon)$ strong solution of (SNS) without concentration with the following uniform bound for the difference:
  
  \[ \sup_{0 \leq t < \infty} (\norm{U_h^\epsilon}_{H^1} + \epsilon^2 \norm{U_3^\epsilon}_{H^1}^2) + \int_0^\infty \big( \norm{\nabla U_h^\epsilon}_{H^1}^2 + \epsilon^2 \norm{\nabla U_3^\epsilon}_{H^1}^2 \big) \diff t \leq C \epsilon^2,\]
  
  where C depends only on the domain size and $\norm{{u_h}_0}_{H^2}.$
  
  Based on the existence of $(\mathbf{u}_h^\epsilon, u_3^\epsilon)$, the existence of the extended $(\mathbf{u}_h^\epsilon, u_3^\epsilon, c^\epsilon)$ strong solution for (E-SNS) follows from identical steps as the existence of the velocity part for the scaled Navier-Stokes equations with fix $\epsilon$ (in the sense that the same ideas can be applied to the extended solution vector as for the original velocity-only solution). However, with this method we only know the local existence of $c^\epsilon.$ (Because in general what we know is the local unique existence of strong solutions, the global existence for the velocity came only as a result of considering small enough $\epsilon$ values). The global existence for small epsilon values of the $(\mathbf{u}_h^\epsilon, u_3^\epsilon)$ solution was a result of Li-Titi's article. Here we will propose the global existence of $c^\epsilon$ later.
  
  \item \textbf{$H^4$ space regularity for $\mathbf{u}_h.$} Based on a comment in Li-Titi's Strong Convergence Justification article, page10: in general when ${u_h}_0$ has more regularities, then correspondingly the solution $u_h$ will have better regularities as well, see Petcu-Wirosoetisno. With ${u_h}_0 \in H^3(\Omega)$ we have a \todo{why the loc here}
  
  \begin{equation} \label{H4_uh}
    u_h \in L^2_{\text{loc}} H^4
  \end{equation}
  
  solution function. 
  
  \item \textbf{$H^3$ space estimate for $U_h^\epsilon.$} Based on Remark 1.1/(iii) in Li-Titi's Strong Convergence Justification article, page6, we have for ${u_h}_0 \in H^k, k \geq 2$ the following general result:
  
  \begin{equation}
    \sup_{t}{\norm{U_h^\epsilon}_{H^{k-1}}^2} + \int_0^\infty \norm{\nabla U_h^\epsilon}_{H^{k-1}}^2 \leq C \epsilon^2.
  \end{equation}
  
  With $k=3$ this provides that the $L^2 H^3$ norm of $U_h^\epsilon$ is bounded by $C \epsilon^2:$
  
  \begin{equation} \label{H3_Uh}
    \int_0^\infty \norm{U_h^\epsilon}_{H^3}^2 \leq C \epsilon^2.
  \end{equation}

\end{enumerate}

\subsection{$c^\epsilon \to c$ in $L^\infty H^1$ (estimating $\nabla C^\epsilon$)}

 It seems that whatever regularity we can need for $C^\epsilon$ here in this coordinate system, it will be based and built on the $\partial_t C^\epsilon$ term, because from the Laplacians we won't have a complete term as the $\epsilon$ multiplier is there.
    
    Multiplying the governing equation

$$ \partial_t C^\epsilon + U^\epsilon \nabla C^\epsilon  + U^\epsilon \nabla c+ u \nabla C^\epsilon = \textcolor{orange}{\epsilon K_1 \partial_{11} C^\epsilon + \epsilon K_1 \partial_{11} c} + K_2 \partial_{22} C^\epsilon + K_3 \partial_{33} C^\epsilon + S^\epsilon $$

by $- \partial_{ii} C^\epsilon, i=1,2,3$ we obtain

\begin{equation}
  \begin{split}
      & \frac{1}{2}\frac{\diff}{\diff t} \norm{\nabla C^\epsilon}_2^2 + \epsilon K_1 \norm{\partial_{11} C^\epsilon}_2^2 + K_2 \norm{\partial_{22} C^\epsilon}_2^2 + K_3 \norm{\partial_{33} C^\epsilon}_2^2 + C \norm{\partial_{ij} C^\epsilon}_2^2 \\
      & = \sum_{i = 1,2,3} \int_\Omega U^\epsilon \nabla C^\epsilon \partial_{ii} C^\epsilon + \int_\Omega U^\epsilon \nabla c \partial_{ii} C^\epsilon + \int_\Omega u \nabla C^\epsilon \partial_{ii} C^\epsilon - \int_\Omega \epsilon K_1 \partial_{11} c \partial_{ii} C^\epsilon \\
      & = \sum_{i = 1,2,3} I_1(i)  + I_2(i) + I_3(i) + I_4(i),
  \end{split}
\end{equation}

where $i \neq j$ and for now we skip the $S^\epsilon$ function.

Now we estimate the terms on the right-hand side one by one.

\begin{equation}
  \begin{split}
    & I_1(1) = \int_\Omega U^\epsilon \nabla C^\epsilon \partial_{11} C^\epsilon \\
    & = \int_\Omega U_1^\epsilon \partial_1 C^\epsilon \partial_{11} C^\epsilon + \int_\Omega U_2^\epsilon \partial_2 C^\epsilon \partial_{11} C^\epsilon + \int_\Omega U_3^\epsilon \partial_3 C^\epsilon \partial_{11} C^\epsilon \\
    & = I_1(1)(a) + I_1(1)(b) + I_1(1)(c).
  \end{split}
\end{equation}

Firstly we estimate $I_1(1)(a).$ We have

$$ \int_\Omega U_1^\epsilon \partial_1 C^\epsilon \partial_{11} C^\epsilon = - \int_\Omega \partial_1 U_1^\epsilon (\partial_1 C^\epsilon)^2 - \int_\Omega U_1^\epsilon \partial_{11} C^\epsilon \partial_1 C^\epsilon,$$

and thus

$$\int_\Omega U_1^\epsilon \partial_1 C^\epsilon \partial_{11} C^\epsilon = - \frac{1}{2} \int_\Omega \partial_1 U_1^\epsilon (\partial_1 C^\epsilon)^2.$$

Using the Holder inequality and $\norm{\cdot}_{L^\infty} \leq \norm{\cdot}_{H^2},$ we have

\begin{equation}
  \begin{split}
    & I_1(1)(a) = \int_\Omega U_1^\epsilon \partial_1 C^\epsilon \partial_{11} C^\epsilon \\
    & \leq C \norm{\partial_1 U_1^\epsilon}_{L^\infty} \cdot \norm{\partial_1 C^\epsilon}_2^2 \leq C \norm{\partial_1 U_1^\epsilon}_{H^2} \cdot \norm{\nabla C^\epsilon}_2^2 \\
    & \leq C \norm{U_1^\epsilon}_{H^3} \cdot \norm{\nabla C^\epsilon}_2^2.
  \end{split}
\end{equation}

The above calculation is one of the points throughout the proof where the $H^3$ regularity is required sharply, otherwise the closure is not possible. 

We can estimate $I_1(1)(b), I_1(1)(c)$ in one common step. For $j = 2,3$ we have

$$ \int_\Omega U_j^\epsilon \partial_j C^\epsilon \partial_{11} C^\epsilon = - \int_\Omega \partial_1 U_j^\epsilon \partial_j C^\epsilon \partial_1 C^\epsilon - \int_\Omega U_j^\epsilon \partial_{j1} C^\epsilon \partial_1 C^\epsilon.$$

Firstly, by the Holder, Ladyzhenskaya and Young inequalities we have
\begin{equation}
  \begin{split}
    & \int_\Omega \partial_1 U_j^\epsilon \partial_j C^\epsilon \partial_1 C^\epsilon \leq \| \partial_1 U_j^\epsilon \| _3 \norm{\partial_j C^\epsilon}_6 \norm{\partial_1 C^\epsilon}_2 \\
    & \leq \| \partial_1 U_j^\epsilon \|_2^{1/2} \|\nabla \partial_1 U_j^\epsilon \|_2^{1/2} \norm{\nabla \partial_j C^\epsilon}_2 \norm{\partial_1 C^\epsilon}_2 \\
    & \leq \zeta \norm{\nabla \partial_j C^\epsilon}_2^2 + \frac{1}{\zeta} \| \partial_1 U_j^\epsilon \|_2 \|\nabla \partial_1 U_j^\epsilon \|_2 \norm{\partial_1 C^\epsilon}_2^2 \\
    & \leq \zeta \norm{\nabla \partial_j C^\epsilon}_2^2 + \frac{1}{\zeta} \|\nabla \partial_1 U_j^\epsilon \|_2^2 \norm{\nabla C^\epsilon}_2^2 \\
    & \leq \zeta \norm{\nabla \partial_j C^\epsilon}_2^2 + \frac{1}{\zeta} \norm{U_h}_{H^3}^2 \norm{\nabla C^\epsilon}_2^2.
  \end{split}
\end{equation}

Note that for $j=3$ this approach -- which is different from what we applied for $I_1(1)(a)$ -- is necessary, as for $\| \partial_1 U_3^\epsilon \| _ {L^\infty} \leq \norm{U_3^\epsilon}_{H^3} \leq \| U_h^\epsilon \| _{H^4}$ we do not have a bound. Since in the $\int_\Omega U_j^\epsilon \partial_{j1} C^\epsilon \partial_1 C^\epsilon$ second term we do not have a derivative on the velocity function, we can simply use

\begin{equation} \nonumber
  \begin{split}
    & \int_\Omega U_j^\epsilon \partial_{j1} C^\epsilon \partial_1 C^\epsilon \leq \| U_j^\epsilon \| _{L^\infty} \norm{\partial_{j1} C^\epsilon}_2 \norm{\partial_1 C^\epsilon}_2 \\
    &  \leq \zeta \norm{\partial_{j1} C^\epsilon}_2^2 +  \frac{1}{\zeta} \| U_j^\epsilon \| _{L^\infty}^2 \norm{\nabla C^\epsilon}_2^2 \leq \zeta \norm{\partial_{j1} C^\epsilon}_2^2 +  \frac{1}{\zeta} \| U_h^\epsilon \| _{H^3}^2 \norm{\nabla C^\epsilon}_2^2. 
  \end{split}
\end{equation}

All together for $j = 2,3$ we have

\begin{equation}
  \begin{split}
    & I_1(1)(b), I_1(1)(c) = \int_\Omega U_j^\epsilon \partial_j C^\epsilon \partial_{11} C^\epsilon \\
    & \leq \zeta \norm{\nabla \partial_j C^\epsilon}_2^2 + \frac{1}{\zeta} \norm{U_h}_{H^3}^2 \norm{\nabla C^\epsilon}_2^2.
  \end{split}
\end{equation}
    
\subsection{The time-dependent estimate for the system}

Using the results from the previous subsection, we have 

\begin{equation}
  \begin{split}
    & \frac{1}{2} \frac{\diff}{\diff t} \norm{\nabla C^\epsilon}_2^2 + \epsilon K_1 \norm{\partial_{11} C^\epsilon}_2^2 + K_2 \norm{\partial_{22} C^\epsilon}_2^2 + K_3 \norm{\partial_{33} C^\epsilon}_2^2 + C \norm{\partial_{ij} C^\epsilon}_2^2 \\
    & \leq C \textcolor{blue}{\norm{U_H^\epsilon}_{H^3}} \norm{\nabla C^\epsilon}_2^2 + C \norm{U_H^\epsilon}_{H^3}^2 \norm{\nabla C^\epsilon}_2^2 \\
    & + C \textcolor{blue}{\norm{u_H^\epsilon}_{H^3}} \norm{\nabla C^\epsilon}_2^2 + C \norm{u_H^\epsilon}_{H^3}^2 \norm{\nabla C^\epsilon}_2^2 \\
    & + \textcolor{green}{C \norm{U_H^\epsilon}_{H^3}^2} + (\norm{\partial_{ij} c}_2^2 + \norm{\nabla \partial_i c}_2^2)\norm{\nabla C^\epsilon}_2^2 \\
    & + \textcolor{green}{C \norm{U_H^\epsilon}_{H^3}^2 \norm{\nabla c}_2^2} \\
    & + \textcolor{green}{\epsilon K_1 \norm{\partial_{11} c}_2^2},
  \end{split}
\end{equation}

where $i \neq j,$ and the appropriate terms have been merged into the left-hand side.

The terms with "the missing square" are due to the $u_1 \partial_1 C^\epsilon \partial_{11} C^\epsilon, U_1^\epsilon \partial_1 C^\epsilon \partial_{11} C^\epsilon$ terms.

The \textcolor{green}{additive terms} will contribute

\[ \int_0^T \epsilon C \norm{\partial_{11} c}_2^2 + C \norm{U_H^\epsilon}_{H^3}^2 \norm{\nabla c}_2^2 + C \norm{U_H^\epsilon}_{H^3}^2  \diff \tau \leq \epsilon C + \epsilon^2 C \]

to the estimate obtained after applying the Gronwall inequality, which is independent from $T$, and the estimate holds even for $T= \infty.$

The multiplicative terms contribute

\[ \int_0^T \textcolor{blue}{\norm{U_H^\epsilon}_{H^3}} + \textcolor{blue}{\norm{u_H^\epsilon}_{H^3}} + \norm{U_H^\epsilon}_{H^3}^2 + C \norm{u_H^\epsilon}_{H^3}^2 + (\norm{\partial_{ij} c}_2^2 + \norm{\nabla \partial_i c}_2^2) \diff \tau,\]

which can be estimated with

\begin{equation}
  \begin{split}
    & \int_0^T \norm{U_H^\epsilon}_{H^3} + \norm{u_H^\epsilon}_{H^3} \diff \tau + C \\
    & \leq \int_0^T \norm{1 \cdot U_H^\epsilon}_{H^3} + \norm{1 \cdot u_H^\epsilon}_{H^3} \diff \tau + C \\
    & \leq T^{1/2} \bigg (\int_0^T \norm{U_H^\epsilon}_{H^3}^2 \bigg )^{1/2} \diff \tau + T^{1/2} \bigg (\int_0^T \norm{u_H^\epsilon}_{H^3}^2 \bigg )^{1/2} \diff \tau + C \\
    & \leq T^{1/2} C + C \\
    & \leq C \sqrt{T}.
  \end{split}
\end{equation}

The issue here is that there is no known available estimate for $L^1$ estimate in time for the Navier-Stokes equations for infinite time interval, only for the case of a bounded domain, where $L^2$ boundedness implies $L^1$ boundedness directly.

Thus the main difference from Section 5 in Li-Titi is that while in their case they had an $\epsilon^2 C$ bound where $C=C(L_1, L_2, \norm{v_0}_{H_1}),$ \textbf{we have a $C$ that depends also on $T.$}

With this we get

\[ \norm{C^\epsilon}_{L^\infty(0,T; H^1(\Omega))} \leq \epsilon C \exp{(C \sqrt{T})}.\]

\textbf{Points:}
\begin{enumerate}
  \item The main question is if we can get $L_1$ estimates in time, which right now I don't know of we can, because applying a square root on the inequality before integrating in time I think might not work because the time derivative is negative. The space norm in itself would contain a square root over it, applying that on the LHS is problematic because of the presence of the time derivative. If we can not get $L_1$ estimates, we will have the $\sqrt{T}$ multiplier term in the Gronwall inequality.
  \item From Proposition 5.3 we know already that the strong solutions $v^\epsilon, w^\epsilon$ of the anisotropic system exist \textbf{globally}. We also know that most likely (still to be checked though) $(v,w,c)$ exist globally, too.
  \item We also do not have the $\delta_0$ part, which they need to deal with. We have a different structure and a higher regularity and we don't arrive to that split structure in that sense. We do not have this $\delta_0$ step. We either have a bounded $C^\epsilon$ norm immediately for any $T$ (if the $L^1$ norm in time turns out to be bounded for the term where the square is missing), then we are fine until infinity, otherwise we don't have a $T$-free estimate, in that case the convergence gets weaker in time.
  \item Which is why, it seems like that the estimate for $c^\epsilon - c$ could work in one piece, taking $T$ which is the maximum existence for $c^\epsilon,$ but if the energy bound works, then $T$ naturally becomes infinite.
  \item With these we can estimate $C^\epsilon$, which is bounded for any finite time and does not blow up at any finite point in time. Thus, I think $c^\epsilon$ also exists for any $T$ smaller than infinity.
  \item Difference: at Li-Titi they have $C,$ we have $C(T)$ unless we can bound the $L_1$ norm in time by a $C$ constant.
  \item We can not use $T=\infty$ in the $C^\epsilon$ bound itself, but the estimate holds for any finite point in time --- the convergence gets slower as time evolves.
  \item If the $L^1$ norm in time is not bounded, then it seems that not just the convergence rate weakens (from $\epsilon^2$ to $\epsilon$), but also the uniformity in time is lost.
\end{enumerate}

\section{Domain structure}

What we need is a scenario where the $H^2$ norm is equivalent to the $L^2$ norm of the Laplacian. There are two options for this:
\begin{itemize}
  \item a periodic domain
  \item a minimal set of constraints would be for a function $f$ to require that $f, \nabla f$ is zero on the boundary.
\end{itemize}

\textbf{We choose the local physical boundary for consistency reasons and set that the velocity and concentration values and the first derivatives of the velocity and concentration are zero.}
This is important because we need to keep the structure of the inequalities we use in Section 5. These inequalities are obtained by using Lemma 2.2 in Li-Titi, which is formulated for periodic domains.

In case we don't have any of the scenarios above, and we can not benefit from the Poincare inequality, then instead of the original

$$ \abs{\int_\Omega (\varphi \cdot \nabla \Phi) \Psi \diff x \diff y \diff z} \leq C \norm{\nabla \varphi_H}_2^{1/2} \norm{\Delta \varphi_H}_2^{1/2} \norm{\nabla \Phi}_2^{1/2} \norm{\Delta \Phi}_2^{1/2} \norm{\Psi}_2 $$

inequality, we have

\begin{equation}
  \begin{split}
    & C \bigg [ \norm{\nabla \varphi_H}_2^{1/2} \big ( \norm{\nabla \varphi_H}_2 + \norm{\nabla_H \varphi_H}_2 + \norm{\Delta_H \varphi_H}_2\big )^{1/2} \\
    & + \norm{\nabla \varphi_H}_2^{1/2} \big ( \norm{\nabla \varphi_H}_2 + \norm{\partial_z \nabla_H \varphi_H}_2 + \norm{\Delta_H \varphi_H}_2 \big )^{1/2} \bigg ] \\
    & \times \norm{\nabla \Phi}_2^{1/2}  \big ( \norm{\nabla \Phi}_2 + \norm{\nabla_H \nabla \Phi}_2 \big )^{1/2} \norm{\Psi}_2.
  \end{split}
\end{equation}

on the right-hand side. This also means that whenever we apply it in this form, it alters the structure significantly. For example, applying it on (5.4) in Li-Titi, the estimate of the term $\int (U \cdot \nabla) V \cdot \Delta V$ will have terms such as $\norm{\nabla V}_2^2 \norm{\nabla_H \nabla V} \norm{\partial_z \nabla_H V},$ which originally fell into the "box" of $\norm{\nabla V_2^2} \norm{\Delta V}_2^2,$ thanks to the Poincare inequality.

This is a significant change, because while the term $\norm{\nabla V_2^2} \norm{\Delta V}_2^2$ can be handled as $\delta_0 \norm{\Delta V}_2^2,$ and as such it can be brought to the left-hand side when creating the structure for the Gronwall inequality, the $$\norm{\nabla V}_2^2 \norm{\nabla_H \nabla V} \norm{\partial_z \nabla_H V}$$ term is problematic because its second half is not present in the left-hand side. Structurally, $ \norm{\nabla_H \nabla V} \norm{\partial_z \nabla_H V}$ does not belong to the left, neither to the right-hand side.

This is why we choose to define conditions where Lemma 2.2 holds in its original form.

\section{Questions}

Questions listed in the previous version of this file.

\end{document}

